% Options for packages loaded elsewhere
\PassOptionsToPackage{unicode}{hyperref}
\PassOptionsToPackage{hyphens}{url}
\documentclass[
  11pt,
  a4paper,
]{article}
\usepackage{xcolor}
\usepackage[margin=18mm]{geometry}
\usepackage{amsmath,amssymb}
\setcounter{secnumdepth}{-\maxdimen} % remove section numbering
\usepackage{iftex}
\ifPDFTeX
  \usepackage[T1]{fontenc}
  \usepackage[utf8]{inputenc}
  \usepackage{textcomp} % provide euro and other symbols
\else % if luatex or xetex
  \usepackage{unicode-math} % this also loads fontspec
  \defaultfontfeatures{Scale=MatchLowercase}
  \defaultfontfeatures[\rmfamily]{Ligatures=TeX,Scale=1}
\fi
\ifPDFTeX\else
  % xetex/luatex font selection
    \setmainfont[]{Noto Serif}
    \setsansfont[]{Noto Sans}
    \setmonofont[]{Noto Sans Mono}
  \setmathfont[]{Noto Sans Math}
\fi
% Use upquote if available, for straight quotes in verbatim environments
\IfFileExists{upquote.sty}{\usepackage{upquote}}{}
\IfFileExists{microtype.sty}{% use microtype if available
  \usepackage[]{microtype}
  \UseMicrotypeSet[protrusion]{basicmath} % disable protrusion for tt fonts
}{}
\makeatletter
\@ifundefined{KOMAClassName}{% if non-KOMA class
  \IfFileExists{parskip.sty}{%
    \usepackage{parskip}
  }{% else
    \setlength{\parindent}{0pt}
    \setlength{\parskip}{6pt plus 2pt minus 1pt}}
}{% if KOMA class
  \KOMAoptions{parskip=half}}
\makeatother
\ifLuaTeX
\usepackage[bidi=basic]{babel}
\else
\usepackage[bidi=default]{babel}
\fi
\babelprovide[main,import]{russian}
\ifPDFTeX
\else
\babelfont{rm}[]{Noto Serif}
\fi
% get rid of language-specific shorthands (see #6817):
\let\LanguageShortHands\languageshorthands
\def\languageshorthands#1{}
\setlength{\emergencystretch}{3em} % prevent overfull lines
\providecommand{\tightlist}{%
  \setlength{\itemsep}{0pt}\setlength{\parskip}{0pt}}
\usepackage{bookmark}
\IfFileExists{xurl.sty}{\usepackage{xurl}}{} % add URL line breaks if available
\urlstyle{same}
\hypersetup{
  pdftitle={Билеты по дискретным случайным процессам},
  pdflang={ru-RU},
  hidelinks,
  pdfcreator={LaTeX via pandoc}}

\title{Билеты по дискретным случайным процессам}
\author{}
\date{}

\begin{document}
\maketitle

\section{Билет 06. Распределения случайных величин, векторов и
элементов. Конечномерные распределения и
плотности}\label{ux431ux438ux43bux435ux442-06.-ux440ux430ux441ux43fux440ux435ux434ux435ux43bux435ux43dux438ux44f-ux441ux43bux443ux447ux430ux439ux43dux44bux445-ux432ux435ux43bux438ux447ux438ux43d-ux432ux435ux43aux442ux43eux440ux43eux432-ux438-ux44dux43bux435ux43cux435ux43dux442ux43eux432.-ux43aux43eux43dux435ux447ux43dux43eux43cux435ux440ux43dux44bux435-ux440ux430ux441ux43fux440ux435ux434ux435ux43bux435ux43dux438ux44f-ux438-ux43fux43bux43eux442ux43dux43eux441ux442ux438}

Источники: Ширяев 1, гл. II; Сердобольская, лекция 1; лекции ДСП.

\subsection{1. Короткий
ответ}\label{ux43aux43eux440ux43eux442ux43aux438ux439-ux43eux442ux432ux435ux442}

Распределение случайного элемента - это образ вероятностной меры при
измеримом отображении. Если

\[
X:(\Omega,\mathcal F,P)\to(E,\mathcal E)
\]

измерим, то его распределение задается формулой

\[
P_X(B)=P\{X\in B\}=P(X^{-1}(B)),\qquad B\in\mathcal E.
\]

Для вещественной случайной величины распределение можно задавать
функцией распределения

\[
F_X(x)=P\{X\le x\}.
\]

Для случайного вектора \(X=(X_1,\ldots,X_d)\):

\[
F_X(x_1,\ldots,x_d)=P\{X_1\le x_1,\ldots,X_d\le x_d\}.
\]

Для процесса \(\{X_t:t\in T\}\) конечномерные распределения - это
распределения векторов

\[
(X_{t_1},\ldots,X_{t_n}).
\]

Если распределение имеет плотность \(p\), то вероятность множества
считается интегралом:

\[
P_X(B)=\int_B p(x)\,dx.
\]

\subsection{2. Полный
ответ}\label{ux43fux43eux43bux43dux44bux439-ux43eux442ux432ux435ux442}

\subsubsection{6.1. Введение: перенос вероятностной
меры}\label{ux432ux432ux435ux434ux435ux43dux438ux435-ux43fux435ux440ux435ux43dux43eux441-ux432ux435ux440ux43eux44fux442ux43dux43eux441ux442ux43dux43eux439-ux43cux435ux440ux44b}

В Билете 05 случайная величина была определена как измеримое
отображение. Теперь главный вопрос: какую вероятность она задает на
пространстве своих значений. Исходная вероятность \(P\) живет на
\((\Omega,\mathcal F)\), но наблюдаем мы обычно не сам исход \(\omega\),
а значение \(X(\omega)\). Поэтому вероятность нужно перенести с
\(\Omega\) на пространство значений.

\subsubsection{6.2. Определение распределения
(закона)}\label{ux43eux43fux440ux435ux434ux435ux43bux435ux43dux438ux435-ux440ux430ux441ux43fux440ux435ux434ux435ux43bux435ux43dux438ux44f-ux437ux430ux43aux43eux43dux430}

Пусть

\[
X:(\Omega,\mathcal F,P)\to(E,\mathcal E)
\]

\begin{itemize}
\tightlist
\item
  случайный элемент. Его распределением, или законом, называется мера
  \(P_X\) на \((E,\mathcal E)\):
\end{itemize}

\[
P_X(B)=P\{\omega:X(\omega)\in B\}=P(X^{-1}(B)),\qquad B\in\mathcal E.
\]

Эта формула корректна именно потому, что \(X\) измерим: для каждого
\(B\in\mathcal E\) прообраз \(X^{-1}(B)\) лежит в \(\mathcal F\), значит
имеет вероятность.

\subsubsection{6.3. Свойства распределения как вероятностной
меры}\label{ux441ux432ux43eux439ux441ux442ux432ux430-ux440ux430ux441ux43fux440ux435ux434ux435ux43bux435ux43dux438ux44f-ux43aux430ux43a-ux432ux435ux440ux43eux44fux442ux43dux43eux441ux442ux43dux43eux439-ux43cux435ux440ux44b}

Распределение является вероятностной мерой. Действительно,

\[
P_X(E)=P(X^{-1}(E))=P(\Omega)=1,
\]

а счетная аддитивность следует из того, что прообраз сохраняет
дизъюнктные объединения:

\[
X^{-1}\left(\bigsqcup_{n=1}^{\infty}B_n\right)
=
\bigsqcup_{n=1}^{\infty}X^{-1}(B_n).
\]

\subsubsection{6.4. Функция распределения вещественной случайной
величины}\label{ux444ux443ux43dux43aux446ux438ux44f-ux440ux430ux441ux43fux440ux435ux434ux435ux43bux435ux43dux438ux44f-ux432ux435ux449ux435ux441ux442ux432ux435ux43dux43dux43eux439-ux441ux43bux443ux447ux430ux439ux43dux43eux439-ux432ux435ux43bux438ux447ux438ux43dux44b}

Для вещественной случайной величины \(X\) распределение \(P_X\) - это
вероятностная мера на \(\mathcal B(\mathbb R)\). Часто ее задают не
самой мерой, а функцией распределения

\[
F_X(x)=P\{X\le x\}=P_X((-\infty,x]).
\]

Функция распределения не убывает, непрерывна справа и имеет пределы

\[
\lim_{x\to-\infty}F_X(x)=0,\qquad
\lim_{x\to+\infty}F_X(x)=1.
\]

Наоборот, всякая функция с этими свойствами задает вероятностную меру на
\(\mathcal B(\mathbb R)\); это уже использовалось в Билете 02.

\subsubsection{6.5. Распределение и функция распределения случайного
вектора}\label{ux440ux430ux441ux43fux440ux435ux434ux435ux43bux435ux43dux438ux435-ux438-ux444ux443ux43dux43aux446ux438ux44f-ux440ux430ux441ux43fux440ux435ux434ux435ux43bux435ux43dux438ux44f-ux441ux43bux443ux447ux430ux439ux43dux43eux433ux43e-ux432ux435ux43aux442ux43eux440ux430}

Для случайного вектора

\[
X=(X_1,\ldots,X_d)
\]

распределение \(P_X\) - это вероятностная мера на
\(\mathcal B(\mathbb R^d)\). Его функция распределения:

\[
F_X(x_1,\ldots,x_d)
=
P\{X_1\le x_1,\ldots,X_d\le x_d\}.
\]

Она задает вероятности прямоугольников через разности значений \(F_X\).
Например, в размерности \(2\):

\[
P\{a<X\le b,\ c<Y\le d\}
=
F(b,d)-F(a,d)-F(b,c)+F(a,c).
\]

В размерности \(d\) используется аналогичная формула
включений-исключений по вершинам прямоугольника.

\subsubsection{6.6. Комплексные случайные
величины}\label{ux43aux43eux43cux43fux43bux435ux43aux441ux43dux44bux435-ux441ux43bux443ux447ux430ux439ux43dux44bux435-ux432ux435ux43bux438ux447ux438ux43dux44b}

Комплексная случайная величина рассматривается как двумерный
вещественный вектор:

\[
Z=X+iY
\quad\Longleftrightarrow\quad
(\operatorname{Re}Z,\operatorname{Im}Z)=(X,Y).
\]

Поэтому ее распределение - мера на \(\mathcal B(\mathbb C)\), что
естественно отождествляется с \(\mathcal B(\mathbb R^2)\). Для
комплексного случайного вектора из \(\mathbb C^m\) получают
распределение в \(\mathbb R^{2m}\). Важно: у комплексных чисел нет
естественного линейного порядка, поэтому обычно говорят не о функции
\(P\{Z\le z\}\), а о распределении как мере или о совместной функции
распределения вещественных и мнимых частей.

\subsubsection{6.7. Распределение в метрических и функциональных
пространствах}\label{ux440ux430ux441ux43fux440ux435ux434ux435ux43bux435ux43dux438ux435-ux432-ux43cux435ux442ux440ux438ux447ux435ux441ux43aux438ux445-ux438-ux444ux443ux43dux43aux446ux438ux43eux43dux430ux43bux44cux43dux44bux445-ux43fux440ux43eux441ux442ux440ux430ux43dux441ux442ux432ux430ux445}

Если \(X\) - случайный элемент метрического пространства \((E,\rho)\),
то его распределение - вероятностная мера на борелевской
\(\sigma\)-алгебре \(\mathcal B(E)\). Например, случайная
последовательность имеет распределение на пространстве
\(\mathbb R^{\mathbb N}\) с произведенной \(\sigma\)-алгеброй, а
случайная траектория - на выбранном пространстве функций, если это
пространство снабжено подходящей \(\sigma\)-алгеброй.

\subsubsection{6.8. Конечномерные распределения случайного
процесса}\label{ux43aux43eux43dux435ux447ux43dux43eux43cux435ux440ux43dux44bux435-ux440ux430ux441ux43fux440ux435ux434ux435ux43bux435ux43dux438ux44f-ux441ux43bux443ux447ux430ux439ux43dux43eux433ux43e-ux43fux440ux43eux446ux435ux441ux441ux430}

Для случайного процесса \(\{X_t:t\in T\}\) центральный объект -
конечномерные распределения. Для любого набора моментов
\(t_1,\ldots,t_n\) рассматривают случайный вектор

\[
(X_{t_1},\ldots,X_{t_n})
\]

и его распределение:

\[
P_{t_1,\ldots,t_n}(B)
=
P\{(X_{t_1},\ldots,X_{t_n})\in B\},
\qquad B\in\mathcal B(\mathbb R^n).
\]

Если процесс вещественный, его конечномерная функция распределения:

\[
F_{t_1,\ldots,t_n}(x_1,\ldots,x_n)
=
P\{X_{t_1}\le x_1,\ldots,X_{t_n}\le x_n\}.
\]

Эти распределения описывают совместное поведение процесса в конечном
числе моментов времени. Для полного задания закона процесса на
пространстве последовательностей нужно согласованное семейство таких
распределений; это будет темой Билета 15 и Билета 16.

\subsubsection{6.9. Распределения с
плотностью}\label{ux440ux430ux441ux43fux440ux435ux434ux435ux43bux435ux43dux438ux44f-ux441-ux43fux43bux43eux442ux43dux43eux441ux442ux44cux44e}

Отдельный важный случай - распределения с плотностью. Если \(X\)
принимает значения в \(\mathbb R^d\) и существует неотрицательная
функция \(p\) такая, что

\[
\int_{\mathbb R^d}p(x)\,dx=1
\]

и

\[
P_X(B)=\int_B p(x)\,dx,
\]

то \(p\) называется плотностью распределения \(X\) относительно меры
Лебега. В программе этого билета отдельно выделен случай непрерывной
плотности: если \(p\) непрерывна, то вероятности прямоугольников и
достаточно хороших областей можно понимать через обычный несобственный
интеграл Римана. Например,

\[
F_X(x)=\int_{-\infty}^{x}p(u)\,du
\]

в одномерном случае и

\[
F_X(x_1,\ldots,x_d)
=
\int_{-\infty}^{x_1}\cdots\int_{-\infty}^{x_d}
p(u_1,\ldots,u_d)\,du_d\cdots du_1
\]

в \(d\)-мерном случае.

\subsection{3. Основные
определения}\label{ux43eux441ux43dux43eux432ux43dux44bux435-ux43eux43fux440ux435ux434ux435ux43bux435ux43dux438ux44f}

\subsubsection{Распределение случайного
элемента}\label{ux440ux430ux441ux43fux440ux435ux434ux435ux43bux435ux43dux438ux435-ux441ux43bux443ux447ux430ux439ux43dux43eux433ux43e-ux44dux43bux435ux43cux435ux43dux442ux430}

Пусть \(X:(\Omega,\mathcal F,P)\to(E,\mathcal E)\) измерим.
Распределение \(X\) - это мера

\[
P_X=P\circ X^{-1}
\]

на \((E,\mathcal E)\), заданная формулой

\[
P_X(B)=P(X^{-1}(B)),\qquad B\in\mathcal E.
\]

\subsubsection{Закон случайной
величины}\label{ux437ux430ux43aux43eux43d-ux441ux43bux443ux447ux430ux439ux43dux43eux439-ux432ux435ux43bux438ux447ux438ux43dux44b}

То же самое, что распределение. Записи

\[
X\sim\mu,\qquad P_X=\mu
\]

означают, что случайная величина \(X\) имеет распределение \(\mu\).

\subsubsection{Функция распределения вещественной случайной
величины}\label{ux444ux443ux43dux43aux446ux438ux44f-ux440ux430ux441ux43fux440ux435ux434ux435ux43bux435ux43dux438ux44f-ux432ux435ux449ux435ux441ux442ux432ux435ux43dux43dux43eux439-ux441ux43bux443ux447ux430ux439ux43dux43eux439-ux432ux435ux43bux438ux447ux438ux43dux44b-1}

\[
F_X(x)=P\{X\le x\}=P_X((-\infty,x]).
\]

\subsubsection{Атом
распределения}\label{ux430ux442ux43eux43c-ux440ux430ux441ux43fux440ux435ux434ux435ux43bux435ux43dux438ux44f}

Точка \(a\), для которой

\[
P\{X=a\}>0.
\]

Скачок функции распределения в точке \(a\) равен массе атома:

\[
P\{X=a\}=F_X(a)-F_X(a-),
\]

где

\[
F_X(a-)=\lim_{x\uparrow a}F_X(x).
\]

\subsubsection{Дискретное
распределение}\label{ux434ux438ux441ux43aux440ux435ux442ux43dux43eux435-ux440ux430ux441ux43fux440ux435ux434ux435ux43bux435ux43dux438ux435}

Распределение, сосредоточенное на конечном или счетном множестве
\(\{x_k\}\):

\[
P\{X=x_k\}=p_k,\qquad p_k\ge0,\qquad \sum_k p_k=1.
\]

\subsubsection{Распределение случайного
вектора}\label{ux440ux430ux441ux43fux440ux435ux434ux435ux43bux435ux43dux438ux435-ux441ux43bux443ux447ux430ux439ux43dux43eux433ux43e-ux432ux435ux43aux442ux43eux440ux430}

Для

\[
X=(X_1,\ldots,X_d)
\]

это мера \(P_X\) на \(\mathcal B(\mathbb R^d)\):

\[
P_X(B)=P\{(X_1,\ldots,X_d)\in B\}.
\]

\subsubsection{Функция распределения случайного
вектора}\label{ux444ux443ux43dux43aux446ux438ux44f-ux440ux430ux441ux43fux440ux435ux434ux435ux43bux435ux43dux438ux44f-ux441ux43bux443ux447ux430ux439ux43dux43eux433ux43e-ux432ux435ux43aux442ux43eux440ux430}

\[
F_X(x_1,\ldots,x_d)
=
P\{X_1\le x_1,\ldots,X_d\le x_d\}.
\]

\subsubsection{Комплексная случайная
величина}\label{ux43aux43eux43cux43fux43bux435ux43aux441ux43dux430ux44f-ux441ux43bux443ux447ux430ux439ux43dux430ux44f-ux432ux435ux43bux438ux447ux438ux43dux430}

Случайная величина со значениями в \(\mathbb C\), понимаемая как
случайный вектор

\[
Z=\operatorname{Re}Z+i\operatorname{Im}Z
\quad\leftrightarrow\quad
(\operatorname{Re}Z,\operatorname{Im}Z)\in\mathbb R^2.
\]

Ее распределение - мера на
\(\mathcal B(\mathbb C)\simeq\mathcal B(\mathbb R^2)\).

\subsubsection{Случайный элемент метрического
пространства}\label{ux441ux43bux443ux447ux430ux439ux43dux44bux439-ux44dux43bux435ux43cux435ux43dux442-ux43cux435ux442ux440ux438ux447ux435ux441ux43aux43eux433ux43e-ux43fux440ux43eux441ux442ux440ux430ux43dux441ux442ux432ux430}

Измеримое отображение

\[
X:(\Omega,\mathcal F,P)\to(E,\mathcal B(E)),
\]

где \((E,\rho)\) - метрическое пространство. Его распределение -
вероятностная мера на \(\mathcal B(E)\).

\subsubsection{Конечномерное распределение
процесса}\label{ux43aux43eux43dux435ux447ux43dux43eux43cux435ux440ux43dux43eux435-ux440ux430ux441ux43fux440ux435ux434ux435ux43bux435ux43dux438ux435-ux43fux440ux43eux446ux435ux441ux441ux430}

Для процесса \(\{X_t:t\in T\}\) и набора \(t_1,\ldots,t_n\) это
распределение случайного вектора

\[
(X_{t_1},\ldots,X_{t_n}).
\]

\subsubsection{Конечномерная функция
распределения}\label{ux43aux43eux43dux435ux447ux43dux43eux43cux435ux440ux43dux430ux44f-ux444ux443ux43dux43aux446ux438ux44f-ux440ux430ux441ux43fux440ux435ux434ux435ux43bux435ux43dux438ux44f}

\[
F_{t_1,\ldots,t_n}(x_1,\ldots,x_n)
=
P\{X_{t_1}\le x_1,\ldots,X_{t_n}\le x_n\}.
\]

\subsubsection{Плотность
распределения}\label{ux43fux43bux43eux442ux43dux43eux441ux442ux44c-ux440ux430ux441ux43fux440ux435ux434ux435ux43bux435ux43dux438ux44f}

Измеримая функция \(p:\mathbb R^d\to\mathbb R\), неотрицательная почти
всюду, такая, что

\[
\int_{\mathbb R^d}p(x)\,dx=1
\]

и для допустимых множеств \(B\) выполнено

\[
P_X(B)=\int_B p(x)\,dx.
\]

В языке меры это означает, что \(P_X\) абсолютно непрерывно относительно
меры Лебега, а \(p\) является производной Радона-Никодима; формально это
будет связано с Билетом 10.

\subsubsection{Маргинальное
распределение}\label{ux43cux430ux440ux433ux438ux43dux430ux43bux44cux43dux43eux435-ux440ux430ux441ux43fux440ux435ux434ux435ux43bux435ux43dux438ux435}

Распределение части координат случайного вектора. Если \(X=(X_1,X_2)\),
то распределение \(X_1\) получается из совместного распределения:

\[
P_{X_1}(B)=P_{(X_1,X_2)}(B\times\mathbb R).
\]

\subsection{4. Основные теоремы и
свойства}\label{ux43eux441ux43dux43eux432ux43dux44bux435-ux442ux435ux43eux440ux435ux43cux44b-ux438-ux441ux432ux43eux439ux441ux442ux432ux430}

\subsubsection{Свойство 1. Распределение случайного элемента является
вероятностной
мерой}\label{ux441ux432ux43eux439ux441ux442ux432ux43e-1.-ux440ux430ux441ux43fux440ux435ux434ux435ux43bux435ux43dux438ux435-ux441ux43bux443ux447ux430ux439ux43dux43eux433ux43e-ux44dux43bux435ux43cux435ux43dux442ux430-ux44fux432ux43bux44fux435ux442ux441ux44f-ux432ux435ux440ux43eux44fux442ux43dux43eux441ux442ux43dux43eux439-ux43cux435ux440ux43eux439}

Пусть \(X:(\Omega,\mathcal F,P)\to(E,\mathcal E)\) измерим. Тогда
\(P_X\) - вероятностная мера на \((E,\mathcal E)\).

\subsubsection{Доказательство}\label{ux434ux43eux43aux430ux437ux430ux442ux435ux43bux44cux441ux442ux432ux43e}

Во-первых,

\[
P_X(E)=P(X^{-1}(E))=P(\Omega)=1.
\]

Во-вторых, если \(B_1,B_2,\ldots\in\mathcal E\) попарно не пересекаются,
то их прообразы тоже попарно не пересекаются, и

\[
X^{-1}\left(\bigsqcup_{n=1}^{\infty}B_n\right)
=
\bigsqcup_{n=1}^{\infty}X^{-1}(B_n).
\]

По счетной аддитивности \(P\):

\[
P_X\left(\bigsqcup_{n=1}^{\infty}B_n\right)
=
\sum_{n=1}^{\infty}P_X(B_n).
\]

Значит \(P_X\) - вероятностная мера. \(\square\)

\subsubsection{Свойство 2. Распределение вещественной случайной величины
задается функцией
распределения}\label{ux441ux432ux43eux439ux441ux442ux432ux43e-2.-ux440ux430ux441ux43fux440ux435ux434ux435ux43bux435ux43dux438ux435-ux432ux435ux449ux435ux441ux442ux432ux435ux43dux43dux43eux439-ux441ux43bux443ux447ux430ux439ux43dux43eux439-ux432ux435ux43bux438ux447ux438ux43dux44b-ux437ux430ux434ux430ux435ux442ux441ux44f-ux444ux443ux43dux43aux446ux438ux435ux439-ux440ux430ux441ux43fux440ux435ux434ux435ux43bux435ux43dux438ux44f}

Если

\[
F_X(x)=P\{X\le x\},
\]

то значения \(F_X\) однозначно задают меру \(P_X\) на
\(\mathcal B(\mathbb R)\).

\subsubsection{Обоснование}\label{ux43eux431ux43eux441ux43dux43eux432ux430ux43dux438ux435}

Для каждого \(x\)

\[
P_X((-\infty,x])=F_X(x).
\]

Лучами \((-\infty,x]\) порождается \(\mathcal B(\mathbb R)\), поэтому
вероятностная мера на \(\mathcal B(\mathbb R)\) однозначно определяется
значениями на этих лучах. Теорема о существовании меры по функции
распределения формулировалась в Билете 02 без доказательства.

\subsubsection{Свойство 3. Свойства функции распределения на
прямой}\label{ux441ux432ux43eux439ux441ux442ux432ux43e-3.-ux441ux432ux43eux439ux441ux442ux432ux430-ux444ux443ux43dux43aux446ux438ux438-ux440ux430ux441ux43fux440ux435ux434ux435ux43bux435ux43dux438ux44f-ux43dux430-ux43fux440ux44fux43cux43eux439}

Функция распределения \(F\) любой вещественной случайной величины
обладает свойствами:

\begin{enumerate}
\def\labelenumi{\arabic{enumi}.}
\tightlist
\item
  \(F\) не убывает;
\item
  \(F\) непрерывна справа;
\item
  \(\lim_{x\to-\infty}F(x)=0\);
\item
  \(\lim_{x\to+\infty}F(x)=1\).
\end{enumerate}

\subsubsection{Доказательство-идея}\label{ux434ux43eux43aux430ux437ux430ux442ux435ux43bux44cux441ux442ux432ux43e-ux438ux434ux435ux44f}

Монотонность следует из вложения лучей:

\[
x\le y\Rightarrow (-\infty,x]\subset(-\infty,y].
\]

Правые пределы и пределы на бесконечностях следуют из непрерывности
вероятностной меры сверху и снизу для монотонных последовательностей
событий. Например, если \(x_n\downarrow x\), то

\[
(-\infty,x_n]\downarrow(-\infty,x],
\]

поэтому

\[
F(x_n)\downarrow F(x).
\]

\subsubsection{Свойство 4. Масса точки равна скачку функции
распределения}\label{ux441ux432ux43eux439ux441ux442ux432ux43e-4.-ux43cux430ux441ux441ux430-ux442ux43eux447ux43aux438-ux440ux430ux432ux43dux430-ux441ux43aux430ux447ux43aux443-ux444ux443ux43dux43aux446ux438ux438-ux440ux430ux441ux43fux440ux435ux434ux435ux43bux435ux43dux438ux44f}

Для вещественной случайной величины

\[
P\{X=a\}=F(a)-F(a-).
\]

\subsubsection{Доказательство}\label{ux434ux43eux43aux430ux437ux430ux442ux435ux43bux44cux441ux442ux432ux43e-1}

События \(\{X\le x\}\) при \(x\uparrow a\) возрастают к событию
\(\{X<a\}\). Поэтому

\[
F(a-)=P\{X<a\}.
\]

Тогда

\[
P\{X=a\}=P\{X\le a\}-P\{X<a\}=F(a)-F(a-).
\]

\subsubsection{Свойство 5. Вероятность промежутка через функцию
распределения}\label{ux441ux432ux43eux439ux441ux442ux432ux43e-5.-ux432ux435ux440ux43eux44fux442ux43dux43eux441ux442ux44c-ux43fux440ux43eux43cux435ux436ux443ux442ux43aux430-ux447ux435ux440ux435ux437-ux444ux443ux43dux43aux446ux438ux44e-ux440ux430ux441ux43fux440ux435ux434ux435ux43bux435ux43dux438ux44f}

Для \(a<b\):

\[
P\{a<X\le b\}=F(b)-F(a).
\]

Также

\[
P\{X\le b\}=F(b),\qquad
P\{X>b\}=1-F(b).
\]

Остальные варианты интервалов требуют учета атомов на концах.

\subsubsection{Свойство 6. Совместная функция распределения задает
распределение
вектора}\label{ux441ux432ux43eux439ux441ux442ux432ux43e-6.-ux441ux43eux432ux43cux435ux441ux442ux43dux430ux44f-ux444ux443ux43dux43aux446ux438ux44f-ux440ux430ux441ux43fux440ux435ux434ux435ux43bux435ux43dux438ux44f-ux437ux430ux434ux430ux435ux442-ux440ux430ux441ux43fux440ux435ux434ux435ux43bux435ux43dux438ux435-ux432ux435ux43aux442ux43eux440ux430}

Для случайного вектора \(X\in\mathbb R^d\) функция

\[
F_X(x)=P\{X_1\le x_1,\ldots,X_d\le x_d\}
\]

задает значения распределения на прямоугольниках вида

\[
(-\infty,x_1]\times\cdots\times(-\infty,x_d],
\]

а такие прямоугольники порождают \(\mathcal B(\mathbb R^d)\). Поэтому
\(F_X\) однозначно задает распределение \(P_X\).

\subsubsection{Свойство 7. Маргинальные распределения получаются из
совместного}\label{ux441ux432ux43eux439ux441ux442ux432ux43e-7.-ux43cux430ux440ux433ux438ux43dux430ux43bux44cux43dux44bux435-ux440ux430ux441ux43fux440ux435ux434ux435ux43bux435ux43dux438ux44f-ux43fux43eux43bux443ux447ux430ux44eux442ux441ux44f-ux438ux437-ux441ux43eux432ux43cux435ux441ux442ux43dux43eux433ux43e}

Если \(X=(X_1,\ldots,X_d)\) и нужно распределение координат с номерами
\(i_1,\ldots,i_k\), то

\[
P_{(X_{i_1},\ldots,X_{i_k})}(B)
=
P_X(\pi^{-1}(B)),
\]

где \(\pi:\mathbb R^d\to\mathbb R^k\) - соответствующая координатная
проекция.

Для функций распределения это соответствует отправлению лишних
аргументов в \(+\infty\). Например:

\[
F_{X_1}(x)=\lim_{y\to+\infty}F_{(X_1,X_2)}(x,y).
\]

\subsubsection{Свойство 8. Плотность неотрицательна и имеет интеграл
1}\label{ux441ux432ux43eux439ux441ux442ux432ux43e-8.-ux43fux43bux43eux442ux43dux43eux441ux442ux44c-ux43dux435ux43eux442ux440ux438ux446ux430ux442ux435ux43bux44cux43dux430-ux438-ux438ux43cux435ux435ux442-ux438ux43dux442ux435ux433ux440ux430ux43b-1}

Если \(p\) - плотность распределения на \(\mathbb R^d\), то

\[
p(x)\ge0\quad\text{почти всюду}
\]

и

\[
\int_{\mathbb R^d}p(x)\,dx=1.
\]

Вероятность области \(B\) равна интегралу плотности по этой области:

\[
P_X(B)=\int_Bp(x)\,dx.
\]

В одномерном непрерывном случае:

\[
F(x)=\int_{-\infty}^{x}p(u)\,du.
\]

\subsubsection{Свойство 9. Плотность совместного распределения дает
маргинальные плотности
интегрированием}\label{ux441ux432ux43eux439ux441ux442ux432ux43e-9.-ux43fux43bux43eux442ux43dux43eux441ux442ux44c-ux441ux43eux432ux43cux435ux441ux442ux43dux43eux433ux43e-ux440ux430ux441ux43fux440ux435ux434ux435ux43bux435ux43dux438ux44f-ux434ux430ux435ux442-ux43cux430ux440ux433ux438ux43dux430ux43bux44cux43dux44bux435-ux43fux43bux43eux442ux43dux43eux441ux442ux438-ux438ux43dux442ux435ux433ux440ux438ux440ux43eux432ux430ux43dux438ux435ux43c}

Если \((X,Y)\) имеет совместную плотность \(p_{X,Y}(x,y)\), то плотность
\(X\) равна

\[
p_X(x)=\int_{-\infty}^{+\infty}p_{X,Y}(x,y)\,dy,
\]

если этот интеграл корректно определен.

В многомерном случае маргинальная плотность получается интегрированием
совместной плотности по лишним координатам.

\subsubsection{Свойство 10. Конечномерные распределения процесса должны
быть
согласованы}\label{ux441ux432ux43eux439ux441ux442ux432ux43e-10.-ux43aux43eux43dux435ux447ux43dux43eux43cux435ux440ux43dux44bux435-ux440ux430ux441ux43fux440ux435ux434ux435ux43bux435ux43dux438ux44f-ux43fux440ux43eux446ux435ux441ux441ux430-ux434ux43eux43bux436ux43dux44b-ux431ux44bux442ux44c-ux441ux43eux433ux43bux430ux441ux43eux432ux430ux43dux44b}

Если конечномерные распределения действительно получены из одного
процесса, то они согласованы при перестановке координат и при
вычеркивании координат. Например,

\[
P_{t_1,t_2}(B_1\times\mathbb R)=P_{t_1}(B_1),
\]

и

\[
P_{t_2,t_1}(B_2\times B_1)=P_{t_1,t_2}(B_1\times B_2).
\]

Полная теорема о том, когда согласованное семейство задает процесс,
относится к Билету 15.

\subsection{5. Примеры}\label{ux43fux440ux438ux43cux435ux440ux44b}

\subsubsection{Пример 1. Распределение
Бернулли}\label{ux43fux440ux438ux43cux435ux440-1.-ux440ux430ux441ux43fux440ux435ux434ux435ux43bux435ux43dux438ux435-ux431ux435ux440ux43dux443ux43bux43bux438}

Пусть

\[
P\{X=1\}=p,\qquad P\{X=0\}=1-p.
\]

Тогда распределение \(X\) сосредоточено на \(\{0,1\}\):

\[
P_X=(1-p)\delta_0+p\delta_1.
\]

Функция распределения:

\[
F_X(x)=
\begin{cases}
0, & x<0,\\
1-p, & 0\le x<1,\\
1, & x\ge1.
\end{cases}
\]

У этой функции скачки в точках \(0\) и \(1\), равные массам
соответствующих атомов.

\subsubsection{\texorpdfstring{Пример 2. Равномерное распределение на
\([0,1]\)}{Пример 2. Равномерное распределение на {[}0,1{]}}}\label{ux43fux440ux438ux43cux435ux440-2.-ux440ux430ux432ux43dux43eux43cux435ux440ux43dux43eux435-ux440ux430ux441ux43fux440ux435ux434ux435ux43bux435ux43dux438ux435-ux43dux430-01}

Если \(X\) равномерна на \([0,1]\), то

\[
p(x)=
\begin{cases}
1, & x\in[0,1],\\
0, & x\notin[0,1].
\end{cases}
\]

Функция распределения:

\[
F_X(x)=
\begin{cases}
0, & x<0,\\
x, & 0\le x\le1,\\
1, & x>1.
\end{cases}
\]

\subsubsection{Пример 3. Случайный вектор с
плотностью}\label{ux43fux440ux438ux43cux435ux440-3.-ux441ux43bux443ux447ux430ux439ux43dux44bux439-ux432ux435ux43aux442ux43eux440-ux441-ux43fux43bux43eux442ux43dux43eux441ux442ux44cux44e}

Пусть \((X,Y)\) имеет плотность \(p(x,y)\). Тогда

\[
P\{(X,Y)\in B\}=\iint_B p(x,y)\,dx\,dy.
\]

Например,

\[
P\{a<X\le b,\ c<Y\le d\}
=
\int_a^b\int_c^d p(x,y)\,dy\,dx.
\]

Маргинальная плотность \(X\):

\[
p_X(x)=\int_{-\infty}^{+\infty}p(x,y)\,dy.
\]

\subsubsection{Пример 4. Комплексная случайная
величина}\label{ux43fux440ux438ux43cux435ux440-4.-ux43aux43eux43cux43fux43bux435ux43aux441ux43dux430ux44f-ux441ux43bux443ux447ux430ux439ux43dux430ux44f-ux432ux435ux43bux438ux447ux438ux43dux430}

Пусть

\[
Z=X+iY.
\]

Распределение \(Z\) - это распределение вектора \((X,Y)\) на плоскости.
Например,

\[
P\{Z\in B\}=P\{(X,Y)\in \widetilde B\},
\]

где \(\widetilde B\subset\mathbb R^2\) - множество, соответствующее
\(B\subset\mathbb C\) при отождествлении \(\mathbb C\) и
\(\mathbb R^2\).

\subsubsection{Пример 5. Конечномерное распределение случайной
последовательности}\label{ux43fux440ux438ux43cux435ux440-5.-ux43aux43eux43dux435ux447ux43dux43eux43cux435ux440ux43dux43eux435-ux440ux430ux441ux43fux440ux435ux434ux435ux43bux435ux43dux438ux435-ux441ux43bux443ux447ux430ux439ux43dux43eux439-ux43fux43eux441ux43bux435ux434ux43eux432ux430ux442ux435ux43bux44cux43dux43eux441ux442ux438}

Пусть \(X_1,X_2,\ldots\) - случайная последовательность. Для индексов
\(1,3,5\) конечномерное распределение:

\[
P_{1,3,5}(B)=P\{(X_1,X_3,X_5)\in B\},
\qquad B\in\mathcal B(\mathbb R^3).
\]

Его функция распределения:

\[
F_{1,3,5}(x_1,x_3,x_5)
=
P\{X_1\le x_1,\ X_3\le x_3,\ X_5\le x_5\}.
\]

\subsubsection{Пример 6. Распределение случайного
элемента}\label{ux43fux440ux438ux43cux435ux440-6.-ux440ux430ux441ux43fux440ux435ux434ux435ux43bux435ux43dux438ux435-ux441ux43bux443ux447ux430ux439ux43dux43eux433ux43e-ux44dux43bux435ux43cux435ux43dux442ux430}

Пусть

\[
X(\omega)=(X_1(\omega),X_2(\omega),\ldots)
\]

\begin{itemize}
\tightlist
\item
  случайная последовательность как элемент \(\mathbb R^{\mathbb N}\).
  Тогда ее распределение - мера на пространстве последовательностей.
  Цилиндрическое событие
\end{itemize}

\[
C=\{x\in\mathbb R^{\mathbb N}:x_1\le a,\ x_2\in B,\ x_5\le c\}
\]

имеет вероятность

\[
P_X(C)=P\{X_1\le a,\ X_2\in B,\ X_5\le c\}.
\]

\subsection{6. Доказательства и
обоснования}\label{ux434ux43eux43aux430ux437ux430ux442ux435ux43bux44cux441ux442ux432ux430-ux438-ux43eux431ux43eux441ux43dux43eux432ux430ux43dux438ux44f}

\textbf{Почему распределение называют образом меры.}

Формула

\[
P_X(B)=P(X^{-1}(B))
\]

переносит меру \(P\) с исходного пространства на пространство значений.
Поэтому \(P_X\) называют образом меры \(P\) при отображении \(X\) и
пишут

\[
P_X=P\circ X^{-1}.
\]

\textbf{Почему одной функции распределения хватает на прямой.}

На прямой борелевская \(\sigma\)-алгебра порождается лучами
\((-\infty,x]\). Функция \(F_X\) задает вероятности всех таких лучей.
Так как вероятностная мера на порождающем классе при стандартных
условиях продолжается единственным образом, \(F_X\) задает все
распределение.

\textbf{Почему для вектора нужны совместные распределения, а не только
маргинальные.}

Маргинальные распределения \(X\) и \(Y\) не определяют совместное
распределение \((X,Y)\). Например, если \(X\) симметрична и принимает
значения \(\pm1\), то пары

\[
(X,X)
\quad\text{и}\quad
(X,-X)
\]

имеют одинаковые одномерные маргинальные распределения, но разные
совместные распределения.

\textbf{Почему плотность не является вероятностью точки.}

Для непрерывного распределения с плотностью

\[
P\{X=a\}=\int_{\{a\}}p(x)\,dx=0.
\]

Значение \(p(a)\) может быть больше \(1\) и не является вероятностью.
Вероятности получаются только после интегрирования плотности по
множеству положительной длины или объема.

\textbf{Почему комплексный случай сводится к вещественному векторному.}

Пространство \(\mathbb C^m\) естественно отождествляется с
\(\mathbb R^{2m}\):

\[
(z_1,\ldots,z_m)
\leftrightarrow
(\operatorname{Re}z_1,\operatorname{Im}z_1,\ldots,\operatorname{Re}z_m,\operatorname{Im}z_m).
\]

Поэтому распределения комплексных величин - это обычные распределения
вещественных векторов четной размерности.

\textbf{Почему конечномерные распределения процесса должны быть
согласованы.}

Если распределение вектора \((X_{t_1},X_{t_2})\) известно, то
распределение одной координаты \(X_{t_1}\) получается из него
автоматически:

\[
P\{X_{t_1}\in B\}=P\{(X_{t_1},X_{t_2})\in B\times\mathbb R\}.
\]

Поэтому нельзя произвольно назначать распределения для разных наборов
времен: они должны давать одинаковые ответы на пересекающихся наборах
координат.

\subsection{7. Что написать на
доске}\label{ux447ux442ux43e-ux43dux430ux43fux438ux441ux430ux442ux44c-ux43dux430-ux434ux43eux441ux43aux435}

\begin{enumerate}
\def\labelenumi{\arabic{enumi}.}
\tightlist
\item
  Распределение как образ меры:
\end{enumerate}

\[
P_X=P\circ X^{-1},
\qquad
P_X(B)=P\{X\in B\}.
\]

\begin{enumerate}
\def\labelenumi{\arabic{enumi}.}
\setcounter{enumi}{1}
\tightlist
\item
  Функция распределения:
\end{enumerate}

\[
F_X(x)=P\{X\le x\}.
\]

\begin{enumerate}
\def\labelenumi{\arabic{enumi}.}
\setcounter{enumi}{2}
\tightlist
\item
  Случайный вектор:
\end{enumerate}

\[
F_X(x_1,\ldots,x_d)
=
P\{X_1\le x_1,\ldots,X_d\le x_d\}.
\]

\begin{enumerate}
\def\labelenumi{\arabic{enumi}.}
\setcounter{enumi}{3}
\tightlist
\item
  Конечномерные распределения процесса:
\end{enumerate}

\[
P_{t_1,\ldots,t_n}(B)
=
P\{(X_{t_1},\ldots,X_{t_n})\in B\}.
\]

\begin{enumerate}
\def\labelenumi{\arabic{enumi}.}
\setcounter{enumi}{4}
\tightlist
\item
  Конечномерная функция распределения:
\end{enumerate}

\[
F_{t_1,\ldots,t_n}(x_1,\ldots,x_n)
=
P\{X_{t_1}\le x_1,\ldots,X_{t_n}\le x_n\}.
\]

\begin{enumerate}
\def\labelenumi{\arabic{enumi}.}
\setcounter{enumi}{5}
\tightlist
\item
  Плотность:
\end{enumerate}

\[
p\ge0,\qquad
\int_{\mathbb R^d}p(x)\,dx=1,\qquad
P_X(B)=\int_Bp(x)\,dx.
\]

\begin{enumerate}
\def\labelenumi{\arabic{enumi}.}
\setcounter{enumi}{6}
\tightlist
\item
  Маргинализация:
\end{enumerate}

\[
P_{X_1}(B)=P_{(X_1,X_2)}(B\times\mathbb R),
\qquad
p_X(x)=\int p_{X,Y}(x,y)\,dy.
\]

\subsection{8. Типичные дополнительные
вопросы}\label{ux442ux438ux43fux438ux447ux43dux44bux435-ux434ux43eux43fux43eux43bux43dux438ux442ux435ux43bux44cux43dux44bux435-ux432ux43eux43fux440ux43eux441ux44b}

\textbf{Вопрос. Что такое распределение случайной величины?}

Это вероятностная мера на пространстве значений:

\[
P_X(B)=P\{X\in B\}.
\]

Для вещественной случайной величины это мера на
\(\mathcal B(\mathbb R)\).

\textbf{Вопрос. Почему формула \(P_X(B)=P\{X\in B\}\) корректна?}

Потому что \(X\) измерима, значит \(\{X\in B\}=X^{-1}(B)\in\mathcal F\)
для каждого \(B\) из \(\sigma\)-алгебры пространства значений.

\textbf{Вопрос. Чем распределение отличается от функции распределения?}

Распределение - это мера \(P_X\). Функция распределения - числовая
функция \(F_X(x)=P\{X\le x\}\), которая в вещественном случае задает эту
меру.

\textbf{Вопрос. Чем распределение отличается от плотности?}

Распределение - мера. Плотность - функция, которую нужно интегрировать,
чтобы получить вероятности:

\[
P_X(B)=\int_Bp(x)\,dx.
\]

Плотность существует не у всякого распределения.

\textbf{Вопрос. Может ли плотность быть больше единицы?}

Да. Плотность не является вероятностью. Например, равномерное
распределение на \([0,\tfrac12]\) имеет плотность \(2\) на этом
интервале.

\textbf{Вопрос. Как получить распределение координаты из совместного
распределения?}

Нужно взять проекцию или отправить остальные координаты в \(\mathbb R\):

\[
P_{X_1}(B)=P_{(X_1,X_2)}(B\times\mathbb R).
\]

\textbf{Вопрос. Задают ли маргинальные распределения совместное?}

Нет. Они не содержат информации о зависимости между координатами. Для
совместного закона нужна совместная мера или совместная функция
распределения.

\textbf{Вопрос. Как понимать распределение комплексной случайной
величины?}

Как распределение двумерного вещественного вектора
\((\operatorname{Re}Z,\operatorname{Im}Z)\) на плоскости.

\textbf{Вопрос. Что такое конечномерные распределения процесса?}

Это распределения всех конечных наборов координат:

\[
(X_{t_1},\ldots,X_{t_n}).
\]

\textbf{Вопрос. Почему конечномерные распределения должны быть
согласованы?}

Потому что распределение меньшего набора координат должно получаться из
распределения большего набора маргинализацией.

\subsection{9. Типичные
ошибки}\label{ux442ux438ux43fux438ux447ux43dux44bux435-ux43eux448ux438ux431ux43aux438}

\begin{enumerate}
\def\labelenumi{\arabic{enumi}.}
\item
  Путать случайную величину \(X\), ее распределение \(P_X\), функцию
  распределения \(F_X\) и плотность \(p\).
\item
  Говорить о плотности у любого распределения. У дискретного
  распределения относительно меры Лебега плотности нет.
\item
  Считать значение плотности вероятностью: \(p(x)\) - не вероятность
  точки.
\item
  Забывать, что комплексная случайная величина задается двумя
  вещественными координатами.
\item
  Думать, что одномерные маргиналы задают совместное распределение.
\item
  Забывать согласованность конечномерных распределений процесса.
\item
  Использовать формулу \(F_X(x)=P\{X\le x\}\) для объектов без порядка,
  например для общего метрического пространства. Там основной объект -
  мера \(P_X\).
\end{enumerate}

\subsection{10. Связи с другими
билетами}\label{ux441ux432ux44fux437ux438-ux441-ux434ux440ux443ux433ux438ux43cux438-ux431ux438ux43bux435ux442ux430ux43cux438}

\textbf{Билет 02.} Функция распределения на прямой задает вероятностную
меру; это основа перехода от \(F\) к \(P_X\).

\textbf{Билет 04.} Распределения вещественных и векторных случайных
величин живут на борелевских \(\sigma\)-алгебрах.

\textbf{Билет 05.} Измеримость случайного элемента нужна, чтобы формула
\(P_X(B)=P(X^{-1}(B))\) была корректной.

\textbf{Билет 07.} Интегралы по плотности и математические ожидания
требуют аппарата интеграла Лебега.

\textbf{Билет 08.} Характеристическая функция является преобразованием
Фурье распределения.

\textbf{Билет 09.} Независимость случайных величин выражается через
факторизацию совместного распределения.

\textbf{Билет 10.} Плотность в общем виде описывается теоремой
Радона-Никодима.

\textbf{Билет 11.} Замена переменной связывает интегралы по
вероятностному пространству и по распределению.

\textbf{Билет 15.} Согласованные конечномерные распределения задают
случайный процесс.

\textbf{Билет 16.} Дискретный случайный процесс можно рассматривать как
случайный элемент пространства последовательностей; его закон задается
распределением на этом пространстве.

\subsection{11. Минимум на
удовлетворительно}\label{ux43cux438ux43dux438ux43cux443ux43c-ux43dux430-ux443ux434ux43eux432ux43bux435ux442ux432ux43eux440ux438ux442ux435ux43bux44cux43dux43e}

Нужно уметь сказать:

\begin{enumerate}
\def\labelenumi{\arabic{enumi}.}
\tightlist
\item
  Распределение случайного элемента:
\end{enumerate}

\[
P_X(B)=P\{X\in B\}.
\]

\begin{enumerate}
\def\labelenumi{\arabic{enumi}.}
\setcounter{enumi}{1}
\tightlist
\item
  Для вещественной случайной величины:
\end{enumerate}

\[
F_X(x)=P\{X\le x\}.
\]

\begin{enumerate}
\def\labelenumi{\arabic{enumi}.}
\setcounter{enumi}{2}
\tightlist
\item
  Для случайного вектора:
\end{enumerate}

\[
F_X(x_1,\ldots,x_d)
=
P\{X_1\le x_1,\ldots,X_d\le x_d\}.
\]

\begin{enumerate}
\def\labelenumi{\arabic{enumi}.}
\setcounter{enumi}{3}
\item
  Для процесса конечномерные распределения - это распределения векторов
  \((X_{t_1},\ldots,X_{t_n})\).
\item
  При плотности:
\end{enumerate}

\[
P_X(B)=\int_Bp(x)\,dx.
\]

\subsection{12. Минимум на
хорошо}\label{ux43cux438ux43dux438ux43cux443ux43c-ux43dux430-ux445ux43eux440ux43eux448ux43e}

Кроме минимума на удовлетворительно, нужно:

\begin{enumerate}
\def\labelenumi{\arabic{enumi}.}
\tightlist
\item
  Объяснить, почему \(P_X\) является вероятностной мерой.
\item
  Назвать свойства функции распределения: монотонность, правую
  непрерывность, пределы \(0\) и \(1\).
\item
  Показать, как вероятности интервалов выражаются через \(F\).
\item
  Объяснить разницу между распределением, функцией распределения и
  плотностью.
\item
  Рассказать, как получить маргинальное распределение из совместного.
\end{enumerate}

\subsection{13. Что добавить для
отлично}\label{ux447ux442ux43e-ux434ux43eux431ux430ux432ux438ux442ux44c-ux434ux43bux44f-ux43eux442ux43bux438ux447ux43dux43e}

Для сильного ответа стоит добавить:

\begin{enumerate}
\def\labelenumi{\arabic{enumi}.}
\tightlist
\item
  Доказательство переноса меры через прообраз.
\item
  Формулу скачка:
\end{enumerate}

\[
P\{X=a\}=F(a)-F(a-).
\]

\begin{enumerate}
\def\labelenumi{\arabic{enumi}.}
\setcounter{enumi}{2}
\tightlist
\item
  Комплексный случай как распределение вещественного вектора из
  действительных и мнимых частей.
\item
  Согласованность конечномерных распределений процесса.
\item
  Предупреждение, что маргинальные распределения не определяют
  совместное.
\item
  Уточнение: непрерывная плотность позволяет считать вероятности через
  интеграл Римана на прямоугольниках и хороших областях, но общий язык
  плотностей - мера Лебега.
\end{enumerate}

\subsection{14. Устная версия ответа на 5
минут}\label{ux443ux441ux442ux43dux430ux44f-ux432ux435ux440ux441ux438ux44f-ux43eux442ux432ux435ux442ux430-ux43dux430-5-ux43cux438ux43dux443ux442}

Распределение случайного элемента - это способ перенести вероятность с
пространства исходов на пространство значений. Пусть

\[
X:(\Omega,\mathcal F,P)\to(E,\mathcal E)
\]

измерим. Тогда его распределение:

\[
P_X(B)=P\{X\in B\}=P(X^{-1}(B)),\qquad B\in\mathcal E.
\]

Это вероятностная мера на \((E,\mathcal E)\). Измеримость нужна, чтобы
\(X^{-1}(B)\) было событием.

Для вещественной случайной величины распределение можно задавать
функцией распределения:

\[
F_X(x)=P\{X\le x\}.
\]

Она не убывает, непрерывна справа, имеет предел \(0\) при
\(x\to-\infty\) и предел \(1\) при \(x\to+\infty\). Вероятность
полуинтервала выражается как

\[
P\{a<X\le b\}=F(b)-F(a).
\]

Масса точки равна скачку:

\[
P\{X=a\}=F(a)-F(a-).
\]

Для случайного вектора \(X=(X_1,\ldots,X_d)\) распределение - мера на
\(\mathcal B(\mathbb R^d)\), а совместная функция распределения:

\[
F_X(x_1,\ldots,x_d)
=
P\{X_1\le x_1,\ldots,X_d\le x_d\}.
\]

Комплексная случайная величина рассматривается как двумерный
вещественный вектор \((\operatorname{Re}Z,\operatorname{Im}Z)\).

Для процесса \(\{X_t:t\in T\}\) конечномерные распределения - это
распределения векторов

\[
(X_{t_1},\ldots,X_{t_n})
\]

для всех конечных наборов времен. Они должны быть согласованы при
перестановке и вычеркивании координат.

Если распределение имеет плотность \(p\), то

\[
p\ge0\text{ почти всюду},\qquad
\int p(x)\,dx=1,\qquad
P_X(B)=\int_Bp(x)\,dx.
\]

В одномерном случае

\[
F(x)=\int_{-\infty}^{x}p(u)\,du.
\]

Важно не путать распределение, функцию распределения и плотность:
распределение - мера, функция распределения - способ задать меру на
прямой, плотность - функция под интегралом и существует не всегда.

\subsection{15. Если осталось 2 минуты перед
экзаменом}\label{ux435ux441ux43bux438-ux43eux441ux442ux430ux43bux43eux441ux44c-2-ux43cux438ux43dux443ux442ux44b-ux43fux435ux440ux435ux434-ux44dux43aux437ux430ux43cux435ux43dux43eux43c}

Запомнить:

\[
P_X=P\circ X^{-1},
\qquad
P_X(B)=P\{X\in B\}.
\]

Для вещественной случайной величины:

\[
F_X(x)=P\{X\le x\}.
\]

Для вектора:

\[
F_X(x_1,\ldots,x_d)
=
P\{X_1\le x_1,\ldots,X_d\le x_d\}.
\]

Для процесса:

\[
F_{t_1,\ldots,t_n}(x_1,\ldots,x_n)
=
P\{X_{t_1}\le x_1,\ldots,X_{t_n}\le x_n\}.
\]

При плотности:

\[
P_X(B)=\int_Bp(x)\,dx,\qquad
p\ge0,\qquad
\int p=1.
\]

Главная ловушка: \(P_X\), \(F_X\) и \(p\) - разные объекты. Плотность
есть не у всякого распределения.

\newpage

\end{document}
